\documentclass[10pt,a4paper]{amsart}
\usepackage[margin=19mm]{geometry}
\usepackage{amsmath,amssymb,mathtools}
\usepackage[scr=rsfs,cal=euler]{mathalfa}
\usepackage{xfrac}
\usepackage[unicode,hidelinks,bookmarks=false]{hyperref}
\newcommand{\ie}{i.e.\ }
\newcommand{\cf}{cf.\ }
\newcommand{\resp}{respectively\ }
\newcommand{\Subsection}{\subsection}
\providecommand{\nobreakdash}{\nobreak\hbox{-}}
\setlength{\emergencystretch}{2em}
\multlinegap0pt
\usepackage{bm}
\usepackage{bbm}
\usepackage{dsfont}
\usepackage{xspace}
\usepackage{cancel}
\usepackage{mathtools}
\usepackage{amsthm}
\makeatletter
\@ifundefined{theorem}{\newtheorem{theorem}{Theorem}}{}
\@ifundefined{lemma}{\newtheorem{lemma}[theorem]{Lemma}}{}
\makeatother
\newcommand{\End}{\operatorname{End}}
\newcommand{\Hom}{\operatorname{Hom}}
\newcommand{\Ind}{\mathrm{Ind}}
\newcommand{\Ll}{\mathfrak{L}}
\newcommand{\Tr}{\operatorname{Tr}}
\title[Abelian varieties genuinely of $\mathrm{GL}\_n$-type]{Abelian varieties genuinely of $\mathrm{GL}\_n$-type}
\author{Francesc Fit\'e, Enric Florit, Xavier Guitart}
\date{Source: arXiv:2412.21183v3 (CC BY 4.0)}
\begin{document}
\maketitle

\noindent\emph{Source and licence.} Abelian varieties genuinely of $\mathrm{GL}_n$-type. Version arXiv:2412.21183v3 (CC BY 4.0), distributed under the Creative Commons Attribution 4.0 International licence, as recorded in the arXiv licence metadata for this exact version and shown on the abstract page https://arxiv.org/abs/2412.21183v3. Full rights notice: licenses/arxiv-2412.21183-CC-BY-4.0.txt.\par\smallskip

\noindent\textit{Editorial scope.} Counted unit: the Lemma whose source label is \texttt{lemma: unique} in the article (source0.tex lines 890--892, source label \texttt{lemma: unique}), delivered together with the article's own complete original proof of it (source0.tex lines 894--921, one \texttt{proof} environment of 28 lines). No mathematics is written, repaired or re-derived: statement and proof are the article's own text line for line. Required context retained uncounted: lemma: realVL (lines 783--788); equation: dimendk (lines 833--835); equation: HlambdaEnd (lines 841--843). Every definition, display and label of those lines is retained unchanged. Omitted from this package and not redistributed: 1--782, 789--832, 836--840, 844--889, 893--893, 922--1611. Nothing inside a retained excerpt is omitted or altered: every retained body is delivered exactly as the article's lines are written, so no omission receipt of any kind accompanies this package. Citations: the retained excerpts carry no citation command at all, so no bibliography entry of the article is transcribed and no reference is redistributed. Reference adaptation: none was needed and none was made; every reference inside the delivered unit resolves inside it, so no unresolved reference or citation is produced. Printed slips kept literal: no printed word, symbol, hypothesis or number of the counted statement or of its proof was altered. Layout and class: the article's own class is \texttt{amsart} with packages including amsmath, amsthm, amsfonts, amssymb, inputenc, latexsym, verbatim, url, textcomp, xy, color, hyperref, tabularx, xypic; the wrapper is the lane's pinned \texttt{amsart} editorial wrapper at 10pt on A4, which restates the theorem-like environments the retained text needs and adds the article's own macro definitions that the retained text needs (\texttt{\textbackslash End}, \texttt{\textbackslash Hom}, \texttt{\textbackslash Ind}, \texttt{\textbackslash Ll}, \texttt{\textbackslash Tr}), with extra wrapper packages (bm, bbm, dsfont, xspace, cancel, mathtools). The delivered numbering is the unit's own. Assets: no retained span contains a figure environment, a graphics-inclusion command, a PDF-page inclusion, a table or a TikZ picture; no image file is copied into this package, the package declares no asset dependency and no image file is redistributed with it. Licence: version arXiv:2412.21183v3 of this article is distributed under the Creative Commons Attribution 4.0 International licence, as recorded in the arXiv licence metadata for this exact version and shown on the abstract page https://arxiv.org/abs/2412.21183v3. A full rights notice is delivered at licenses/arxiv-2412.21183-CC-BY-4.0.txt. Characters: every retained excerpt is pure ASCII, so the delivered engine, XeLaTeX with its default Latin Modern text font, reports no missing character.
\begin{lemma}\label{lemma: realVL} For every $\lambda \in \Sigma_A$, there is an isomorphism
$$
W_\lambda(A) \otimes_{H_\lambda}E_{\Ll}\simeq V_{\Ll}(A)
$$
 of $E_{\Ll}[G_k]$-modules. In particular, $W_\lambda(A)$ is absolutely irreducible as an $H_\lambda[G_K]$-module for every finite extension $K/k$.
\end{lemma}

\begin{equation}\label{equation: dimendk}
\dim_{H_\lambda}(\End_{H_\lambda[G_k]}(V_\lambda(A)))=[E:H]^2.
\end{equation}

\begin{equation}\label{equation: HlambdaEnd}
\dim_{H_\lambda}(\End_{H_\lambda[G_K]}(V_\lambda(A)))=[E:H]^2.
\end{equation}

\begin{lemma}\label{lemma: unique}
Suppose that $A$ is geometrically of the first kind. Let $\lambda,\lambda'\in \Sigma_A$. There exists at most one finite image character $\chi \colon G_k\rightarrow H_\lambda^ \times$ such that $W_\lambda(A)\simeq \chi\otimes W_{\lambda'} (A)$. In particular, if $\chi$ satisfies that $W_\lambda(A) \simeq \chi \otimes W_\lambda (A)$, then $\chi$ is trivial.
\end{lemma}

\begin{proof}
Let $\chi,\varphi \colon G_k\rightarrow H_\lambda^ \times$ be finite image characters such that 
$$
W_\lambda(A)\simeq \chi\otimes W_{\lambda'} (A)\quad \text{and} \quad W_\lambda(A)\simeq \varphi\otimes W_{\lambda'} (A)\,.
$$
Let $L/k$ be an extension such that $G_L$ is contained in the kernels of $\chi$ and $\varphi$. Lemma \ref{lemma: realVL} implies then that  
$$
\Hom_{G_L}(W_\lambda(A),W_{\lambda'}(A))
$$
has $H_\lambda$-dimension $1$. Since the choice of $L/k$ implies that this representation affords the characters $\chi$ and $\varphi$, we deduce that $\chi$ and $\varphi$ must be equal. 

The second statement of the lemma follows from the first. Let us give an additional slightly different proof.
Let $L/k$ be the field extension cut out by $\chi$. From the isomorphism $W_\lambda(A) \simeq \chi \otimes W_\lambda (A)$, we see that $\Tr(s|V_\lambda(A))=0$ for every $s\in G_k\setminus G_L$.
Hence
\begin{equation}\label{equation: induction formula}
\Ind^{G_L}_{G_k}\big(V_{\lambda}(A)|_{G_L}\big)=\bigoplus_{i=1}^{[L:k]}V_{\lambda}(A).
\end{equation}
By consecutively using \eqref{equation: HlambdaEnd}, Frobenius reciprocity, \eqref{equation: induction formula}, and \eqref{equation: dimendk}, we find
$$
\begin{array}{l@{\,=\,}l}
[E:H]^2 & \displaystyle{ \dim_{H_\lambda} \End_{G_L}(V_\lambda(A))}\\[6pt]
 & \displaystyle{\dim_{H_\lambda} \Hom_{G_k}(V_\lambda(A),\Ind^{G_L}_{G_k}(V_\lambda(A)|_{G_L}))}\\[6pt]
 & \displaystyle{[L:k]\cdot \dim_{H_\lambda} \End_{G_k}(V_\lambda(A))}\\[6pt]
  & \displaystyle{[L:k]\cdot [E:H]^2.}\\[6pt]
\end{array}
$$
Hence $[L:k]=1$ and $\chi$ is trivial.
\end{proof}

\end{document}
